% Image credits for the photographs and AI illustrations of Book 4 (University, Year 2), Arabic edition.
% Machine-readable license records live in images/book4/CREDITS.md; keep this
% file in sync with frontmatter/image-credits-book4.tex (one \item per item there).
% As in the Book 3 Arabic credits: personal names are transliterated, with the
% original spelling in parentheses (the credit itself); institutions, datasets,
% works and licences stay verbatim inside \babelsublr; only the gloss is translated.
\chapter*{مصادر الصور}

الرسوم والخرائط والمخططات في هذا الكتاب أصلية؛ وما رُسم منها من بيانات
منشورة يذكر مصدره في التعليق المرافق أو في هذه القائمة. أما الصور
الفوتوغرافية فمستعملة بمقتضى رخصها الحرة، مع الشكر لأصحابها؛ وروابط المصادر
الكاملة وروابط الرخص لكل صورة مدرجة في \texttt{images/book4/CREDITS.md} في
مستودع الكتاب.

\bigskip
وإلى جانب الصور الفوتوغرافية والرسوم الأصلية، بعض الأشكال رسوم توضيحية
مولَّدة بالذكاء الاصطناعي، أُنتجت بأداة توليد الصور من OpenAI انطلاقًا من
أوامر كتبها محررو الكتاب، ورُوجعت دقتها العلمية قبل إدراجها. وهي لا تُظهر
إلا مشاهد عامة، ولا تمثّل أبدًا مكانًا حقيقيًا أو شخصًا حقيقيًا أو عيّنة
حقيقية. والأمر الكامل لكل صورة مولَّدة مسجَّل في
\texttt{images/book4/ai/PROMPTS.md} في المستودع.

\bigskip
\noindent الصور الفوتوغرافية والبيانات، فصلًا فصلًا:
\begin{itemize}\small
  \item الفصل~1 (الوسط بين النجمي وتكوّن النجوم): القرص الكوكبي الأوّلي حول النجم \babelsublr{HL Tauri}، \babelsublr{ALMA (ESO/NAOJ/NRAO)}، \babelsublr{CC~BY~4.0}؛ وأعمدة الخلق في سديم النسر، \babelsublr{NASA, ESA and the Hubble Heritage Team (STScI/AURA)}، ملك عام.
  \item الفصل~2 (تكوّن الكواكب والنيازك): الكوندريت الكربوني أليندي ونقش ويدمانشتاتن في النيزك الحديدي غيبيون، جيمس سانت جون (\babelsublr{James St.~John})، \babelsublr{CC~BY~2.0}.
  \item الفصل~3 (العناصر في الأرض): دخيل وشاحي من اللرزوليت في البازلت، \babelsublr{Kilbourne Hole}، جيمس سانت جون (\babelsublr{James St.~John})، \babelsublr{CC~BY~2.0}.
  \item الفصل~5 (الصهارات وتموضعها): قاطع قرب شيبروك، دكستر بيركنز (\babelsublr{Dexter Perkins})، \babelsublr{CC0} (عبر \babelsublr{Wikimedia Commons}).
  \item الفصل~6 (الثورانات البركانية): نافورة لابة على مخروط \babelsublr{Pu`u~`\=O`\=o} في كيلاويا، ج.~د.~غريغز (\babelsublr{J.\,D. Griggs})، \babelsublr{USGS}، ملك عام؛ والعمود الثوراني لجبل سانت هيلينز، 18 مايو 1980، أوستن بوست (\babelsublr{Austin Post})، \babelsublr{USGS}، ملك عام؛ والعمود الثوراني لجبل بيناتوبو، 12 يونيو 1991، ديف هارلو (\babelsublr{Dave Harlow})، \babelsublr{USGS}، ملك عام.
  \item الفصل~7 (انصهار الوشاح): عيّنة يدوية من الدونيت، مايك ديفيس (\babelsublr{Mike Davis})، \babelsublr{GeoDIL}، \babelsublr{CC0}.
  \item الفصل~8 (التمايز الصهاري): طبقات الكروميتيت في معقّد بوشفيلد، كيفن والش (\babelsublr{Kevin Walsh})، \babelsublr{CC~BY~2.0}؛ والغابروهات الطبقية في متداخل سكيرغارد، ك.~تيغنر (\babelsublr{C. Tegner})، \babelsublr{CC~BY~4.0}.
  \item الفصل~10 (التفاعلات التحوّلية والسحنات): الشيست الأزرق اللاوسونيتي، \babelsublr{Ward Creek}، كاليفورنيا، والإكلوجيت، ألمينينغن، النرويج، جيمس سانت جون (\babelsublr{James St.~John})، \babelsublr{CC~BY~2.0}.
  \item الفصل~11 (القياس الحراري الضغطي ومسارات الضغط--الحرارة--الزمن): شيست غارنتي، جيمس سانت جون (\babelsublr{James St.~John})، \babelsublr{CC~BY~2.0}؛ وحبيبة كوسيت في الإكلوجيت، بين مستقطِبين متقاطعين، ج.~ر.~سميث (\babelsublr{J.\,R. Smyth})، \babelsublr{CC~BY-SA~1.0}.
  \item الفصل~12 (انصهار القشرة: الميغماتيت والغرانيت): مكشف ميغماتيت في ناشوا، نيوهامبشر، كيفن م.~غيل (\babelsublr{Kevin M. Gill})، \babelsublr{CC~BY~4.0}؛ وعرق من اللوكوغرانيت، آن بيرجس (\babelsublr{Anne Burgess (Geograph)})، \babelsublr{CC~BY-SA~2.0}.
  \item الفصل~13 (تشريح حزام جبلي: جبال الألب): اللابات الوسائدية في أوفيوليت شينايّيه، ليونيل دوشوازيل (\babelsublr{Lionel Duchoiselle})، \babelsublr{CC~BY-SA~3.0}؛ وجبل الماترهورن من زيرمات، \babelsublr{Ximonic (Simo Räsänen)}، \babelsublr{CC~BY-SA~4.0}.
  \item الفصل~14 (مقارنة بين حزام تصادمي وحزام أنديزي): هضبة التبت بين نيبال وشيغاتسي، يان كرانندونك (\babelsublr{Jan Kranendonk})، \babelsublr{CC~BY-SA~4.0}؛ والألتيبلانو البوليفي، لوكا غالوتسي (\babelsublr{Luca Galuzzi})، \babelsublr{CC~BY-SA~2.5}.
  \item الفصل~15 (الكتل القديمة ودورة ويلسون): \babelsublr{Siccar Point}، الجرف الجنوبي الشرقي، \babelsublr{dave souza}، \babelsublr{CC~BY-SA~4.0}.
  \item الفصل~16 (مقارنة بين الأحواض الرسوبية): منصّة حفر شبه غاطسة في بحر الشمال، إريك كريستنسن (\babelsublr{Erik Christensen})، \babelsublr{CC~BY-SA~3.0}.
  \item الفصل~17 (الجزر المحيطية): جزيرة كيور المرجانية من محطة الفضاء الدولية، ناسا (\babelsublr{NASA})، ملك عام.
  \item الفصل~18 (الأرض فيما قبل الكمبري): كوماتييت ذو نسيج السبينيفكس عند \babelsublr{Pyke Hill}، أونتاريو، جيمس سانت جون (\babelsublr{James St.~John})، \babelsublr{CC~BY~2.0}؛ والكربونات الغطائية لتكوين نوكالينا، جنوب أستراليا، جيمس سانت جون (\babelsublr{James St.~John})، \babelsublr{CC~BY~2.0}.
  \item الفصل~19 (الأرض في دهر الحياة الظاهرة): أعمدة الطباشير في \babelsublr{Monument Rocks}، كانساس، جيمس سانت جون (\babelsublr{James St.~John})، \babelsublr{CC~BY~2.0}؛ ومستوى بوناريلّي، \babelsublr{Picenus}، \babelsublr{CC~BY-SA~4.0}.
  \item الفصل~20 (الحفريات والبيئات القديمة): \babelsublr{\emph{Waptia}} و\babelsublr{\emph{Marrella}} من طَفل بيرجس، جيمس سانت جون (\babelsublr{James St.~John})، \babelsublr{CC~BY~2.0}؛ وعيّنة برلين من \babelsublr{\emph{Archaeopteryx}}، ه.~راب (\babelsublr{H.~Raab})، \babelsublr{CC~BY-SA~3.0}؛ و\babelsublr{\emph{Zoophycos}} في ساغريش، \babelsublr{Amada44}، \babelsublr{CC~BY~4.0}.
  \item الفصل~21 (الديناميكا الحرارية للغلاف الجوي): سندان عاصفة رعدية فائقة الخلية، ستيفن كورفيدي (\babelsublr{Stephen Corfidi})، \babelsublr{NOAA}، ملك عام؛ وسبر المسابير اللاسلكية من \babelsublr{Integrated Global Radiosonde Archive}، إ.~دوري وآخرون (\babelsublr{I. Durre and others})، \babelsublr{NOAA NCEI}، ملك عام.
  \item الفصل~22 (الأنظمة الجوية): منخفض جوي فوق أيسلندا، \babelsublr{NASA/GSFC, MODIS Rapid Response Team}، جاك ديكلواتر (\babelsublr{Jacques Descloitres})، ملك عام؛ وعين الإعصار إيزابيل من محطة الفضاء الدولية، إد لو (\babelsublr{Ed Lu})، \babelsublr{NASA}، ملك عام.
  \item الفصل~23 (التقلّب المناخي): مؤشّر تذبذب شمال الأطلسي والمؤشّر المحيطي لإل نينيو، \babelsublr{NOAA Climate Prediction Center}، ملك عام.
  \item الفصل~24 (الأرشيفات والمؤشّرات المناخية البديلة): الطبقات السنوية في لبّ الجليد \babelsublr{GISP2} على عمق 1855~م، أنتوني غاو (\babelsublr{Anthony Gow})، \babelsublr{U.S. Army Corps of Engineers CRREL (NOAA Paleoclimatology)}، ملك عام؛ ومقطع عرضي في رواسب كهفية من الترافرتين، \babelsublr{Diamond Caverns}، كنتاكي، جيمس سانت جون (\babelsublr{James St.~John})، \babelsublr{CC~BY~2.0}؛ والمكدّس القاعي \babelsublr{LR04} لليسيتسكي ورايمو (\babelsublr{L.\,E. Lisiecki and M.\,E. Raymo}، 2005)، وسجلّ حرارة \babelsublr{EPICA Dome C} لجوزيل وآخرين (\babelsublr{J. Jouzel and others}، 2007)، والسجلّ المركَّب لغاز \ce{CO2} في القارة القطبية الجنوبية لبيريتر وآخرين (\babelsublr{B. Bereiter and others}، 2015)، \babelsublr{NOAA NCEI Paleoclimatology}، ملك عام.
  \item الفصل~25 (التأثير المداري والعصور الجليدية): العناصر المدارية والإشماس لبيرجيه ولوتر (\babelsublr{Berger and Loutre}، 1991)، \babelsublr{NOAA NCEI Paleoclimatology}، ملك عام.
  \item الفصل~26 (المناخ عبر الزمن الجيولوجي): السجلّ النظائري القاعي للحقب الحديث \babelsublr{CENOGRID}، فيسترهولد وآخرون (\babelsublr{T. Westerhold and others}، 2020)، \babelsublr{PANGAEA}، \babelsublr{CC~BY~4.0}؛ وتجميع \ce{CO2} في دهر الحياة الظاهرة لفوستر ورويير ولنت (\babelsublr{G.\,L. Foster, D.\,L. Royer and D.\,J. Lunt}، 2017)، \babelsublr{\emph{Nature Communications}}، \babelsublr{CC~BY~4.0}.
  \item الفصل~27 (تغيّر المناخ الناجم عن الإنسان): البيانات الموسَّعة للملحق الثالث من تقرير التقييم السادس للهيئة الحكومية الدولية المعنية بتغيّر المناخ، الفريق العامل الأول، \babelsublr{IPCC AR6 WGI Annex~III}، \babelsublr{CC~BY~4.0}.
  \item الفصل~28 (الأخطار الجيولوجية ومخاطرها): رواسب تدفّق الحطام في كارابالّيدا، فنزويلا، لوسون سميث (\babelsublr{Lawson Smith})، \babelsublr{U.S. Army Corps of Engineers}، ملك عام؛ ومسجد بقي قائمًا بعد تسونامي عام 2004 في آتشيه، \babelsublr{USGS}، ملك عام.
  \item الفصل~29 (مكامن الخامات): منجم بينغهام كانيون من محطة الفضاء الدولية، ناسا (\babelsublr{NASA}، طاقم البعثة~15)، ملك عام؛ والقبّعة الحديدية لمنجم فلامبو، ويسكونسن، جيمس سانت جون (\babelsublr{James St.~John})، \babelsublr{CC~BY~2.0}؛ وبيانات العيار والكتلة الطنّية لسينغر وبيرغر ومورينغ (\babelsublr{D.\,A. Singer, V.\,I. Berger and B.\,C. Moring}، 2008)، \babelsublr{USGS}، ملك عام.
  \item الفصل~30 (جيولوجيا الطاقة): محطّة طاقة حرارية أرضية قرب غيزرفيل، كاليفورنيا، وزارة الدفاع الأمريكية (\babelsublr{U.S. Department of Defense})، ملك عام؛ وكثافات ثنائي أكسيد الكربون من \babelsublr{NIST Chemistry WebBook (Standard Reference Database 69)}.
  \item الفصل~31 (طرائق الحقل والمختبر II): أوجيت بويكيليتي يحيط بالبلاجيوكلاز، \babelsublr{Strekeisen}، \babelsublr{CC~BY-SA~4.0}؛ وشريحة رقيقة من شيست الغارنت والميكا، \babelsublr{Jackdann88}، \babelsublr{CC~BY~3.0}؛ ومسبار إلكتروني دقيق في هيئة المسح الجيولوجي الفنلندية، \babelsublr{Kallerna}، \babelsublr{CC~BY-SA~4.0}؛ ونتائج السيليكا للمواد المرجعية الصخرية، تاغارت وسيمز (\babelsublr{J.\,E. Taggart and D.\,F. Siems}، 2002)، \babelsublr{USGS}، ملك عام.
\end{itemize}
